% !TeX program = xelatex
\documentclass[11pt]{beamer}
\usetheme{nankai}
\usepackage{xeCJK,amsmath,amssymb,booktabs,tabularx,pgfplots}
\setsansfont{Calibri}
\setmainfont{Calibri}
\setCJKmainfont{Microsoft YaHei}[BoldFont=Microsoft YaHei Bold]
\setCJKsansfont{Microsoft YaHei}[BoldFont=Microsoft YaHei Bold]
\usefonttheme[onlymath]{serif}
\pgfplotsset{compat=1.18}
\newcolumntype{Y}{>{\raggedright\arraybackslash}X}
\newcommand{\smalltext}{\fontsize{9pt}{12pt}\selectfont}
\setbeamerfont{frametitle}{size=\large,series=\bfseries}
\setbeamerfont{normal text}{size=\normalsize}
\setbeamersize{text margin left=7mm,text margin right=7mm}
\AtBeginSection[]{}
% 本次汇报不分章，省略模板自动生成的章节副标题。
\setbeamertemplate{frametitle}{%
  \vspace*{-3.5ex}
  \begin{beamercolorbox}[leftskip=2cm]{frametitle}
    \usebeamerfont{frametitle}\insertframetitle\\
    \usebeamerfont{framesubtitle}\strut
  \end{beamercolorbox}
}
\titlebackground*{assets/background}
\title{3×3 2048自主游戏}
\subtitle{强化学习课程实验汇报}
\author{蔡明辰、潘亮、陈达、孟维祎}
\date{}
\footlinepayoff{}
\setbeameroption{hide notes}

\newcommand{\board}[3]{\begingroup\renewcommand{\arraystretch}{1.45}\setlength{\tabcolsep}{9pt}\begin{tabular}{|c|c|c|}\hline#1\\\hline#2\\\hline#3\\\hline\end{tabular}\endgroup}
\newcommand{\ptext}[1]{{\color{revisionpurple}\textbf{#1}}}
\definecolor{revisionpurple}{HTML}{7E0C6E}

\begin{document}
\maketitle
\note{大家好，我们介绍3乘3的2048自主游戏，重点比较搜索方法和强化学习方法的实现与结果。}

\begin{frame}{任务与环境}
\textbf{目标：终局总和尽量大，并生成至少256的方块}\\[2mm]
\smalltext 达到256后继续游戏。每次有效移动后，以0.9/0.1概率落下2/4。\\[3mm]
\begin{columns}[T,onlytextwidth]
\begin{column}{.31\textwidth}\centering 移动前\\[1mm]\board{2&2&}{4&&}{&&}\\[1mm]总和8\end{column}
\begin{column}{.31\textwidth}\centering 向左合并后\\[1mm]\board{4&&}{4&&}{&&}\\[1mm]总和8\end{column}
\begin{column}{.31\textwidth}\centering 随机落块示例\\[1mm]\board{4&&}{4&&}{&&\textcolor{revisionpurple}{2}}\\[1mm]总和10\end{column}
\end{columns}\vspace{3mm}
\smalltext 合并保持总和，随机落块增加总和。释放空间才能继续接纳新方块。\\[1.5mm]
相同的20个测试种子，冻结参数，运行至自然终局。
\end{frame}
\note{目标是让终局棋盘总和尽量大，同时生成至少256的方块。看下面这个例子：两个2合成4，总和仍是8，随机落下一个2后，总和才变成10。所以提高总和，需要不断释放空间来接纳新方块。测试时，达到256后继续游戏，直到没有合法动作，各方法使用同一组20个测试种子。}

\begin{frame}{完整Expectimax：先看完整搜索}
\smalltext 玩家节点选最佳动作，机会节点按真实概率加权，一直递推到终局。
{\large \[Q(s,a)=\sum_{s'}P(s'\mid s,a)V(s'),\qquad V(s)=\max_{a\in A(s)}Q(s,a)
\]}
\begin{columns}[T,onlytextwidth]
\begin{column}{.54\textwidth}
\smalltext $s$为当前棋盘，$s'$为动作与落块后的棋盘。\\[1mm]
$A(s)$为合法动作集合，$P$为转移概率。\\[2mm]
终局边界：$V(s)=U(s)$，$S(s)$为总和
{\large \[U(s)=S(s)+256\cdot\mathbf{1}_{\{\max(s)\ge256\}}
\]}
\end{column}
\begin{column}{.42\textwidth}\smalltext 移动后有2个空格时\\[1.5mm]
\begin{tabular}{@{}lrr@{}}\toprule 位置 & 出现2 & 出现4\\\midrule 空格1 & 0.45 & 0.05\\空格2 & 0.45 & 0.05\\\bottomrule\end{tabular}\\[1.5mm]
各分支概率相加为1。
\end{column}\end{columns}\vspace{1mm}
\smalltext C++求解器，旋转/镜像归一与缓存。首次求解及存储开销较大。
\end{frame}
\note{先讲完整Expectimax，它把规则已知的搜索问题算到终局。在玩家节点，选择期望效用最大的合法动作；在机会节点，按实际落块概率加权。比如移动后有两个空格，每个位置出现2的概率是零点四五，出现4的概率是零点零五。终局效用是棋盘总和，再加256的达标奖励。这里的最优针对这项期望效用，不是优先最大化达标概率。实现采用C加加求解器，用旋转、镜像归一和缓存减少重复局面，代价是首次求解和存储开销较大。}

\begin{frame}{MCTS：用采样减少搜索开销}
\textbf{每次决策128次模拟，搜索深度12，探索系数70}\\[2mm]
\smalltext 未访问动作优先扩展。已访问动作按UCT选择：
{\large \[\mathrm{UCT}(s,a)=\underbrace{\bar Q(s,a)}_{\text{模拟均值}}+\underbrace{c\sqrt{\frac{\ln\!\left(N(s)+1\right)}{N(s,a)}}}_{\text{探索项}}
\]}
\begin{columns}[T,onlytextwidth]
\begin{column}{.48\textwidth}\smalltext
$\bar Q$：该动作的模拟回报均值\\[1.5mm]
$N(s)$：父节点访问次数\\[1.5mm]
$N(s,a)$：动作访问次数，$c$控制探索
\end{column}
\begin{column}{.48\textwidth}\smalltext
每步重新建树，按真实概率模拟落块\\[1.5mm]
到达深度上限后，用启发式评价局面\\[1.5mm]
最终选择根节点平均回报最大的动作
\end{column}\end{columns}\vspace{3mm}
\smalltext 模拟预算决定估计精度，截断估值会影响决策。
\end{frame}
\note{完整枚举的开销较大，MCTS改为采样。每次实际决策重新建树，做128次模拟，最多前看12步。未访问动作优先扩展；已访问动作使用下面的UCT公式。左边是模拟回报的均值，右边是探索项，访问较少的动作得到更大的探索奖励。代码里对父节点访问次数使用加一再取对数。模拟按真实概率落块，深度截断后用启发式估值，最后选择根节点平均回报最大的动作。}

\begin{frame}{TD：从后状态学习长期价值}
\smalltext 后状态$x_t$：移动后、随机落块前的棋盘。以乐观TD为例：
{\large \[\delta_t=\underbrace{z_{t+1}+\gamma V(x_{t+1})}_{\text{更新目标}}-\underbrace{V(x_t)}_{\text{当前估值}}
\]}
\smalltext $z_{t+1}$为新落块数值2或4，$\gamma=1$，终局之后的后续价值取0。\\[1mm]
用TD误差修正当前估值，学习未来可接纳的新方块总和。\\[2mm]
\renewcommand{\arraystretch}{1.04}
\begin{tabularx}{\textwidth}{@{}lY@{}}\toprule
价值表示 & 实现特点\\\midrule
乐观5-tuple & 3张$16^5$表，24个对称特征，初始总价值300\\
CNN-TD & 11通道棋盘编码，卷积价值网络\\
普通n-tuple & 行、列和$2\times2$局部模式，含启发式先验\\\bottomrule
\end{tabularx}\\[2mm]
\smalltext 乐观初始化鼓励尝试较少访问的局面。其余TD版本的效用定义有所不同。
\end{frame}
\note{TD通过训练复用过去的经验。这里用移动后、随机落块前的后状态表示动作结果。以乐观TD为例，价值估计未来还能接纳多少新方块。公式把下一次落块的数值，加上下一后状态的价值，作为更新目标，再减去当前估值。终局之后的后续价值取零，折扣因子取一。乐观5元组将局部棋盘编码成查表特征，用三张表共享24个对称特征，初始总价值设为300。较少访问的局面保持较高估值，能够鼓励探索。我们还比较了CNN和普通元组表，但它们的效用定义有所不同。}

\begin{frame}{PPO：直接学习动作策略}
\textbf{卷积残差网络提取棋盘特征}\\[3mm]
\begin{columns}[T,onlytextwidth]
\begin{column}{.46\textwidth}
\ptext{策略分支}\\[2mm]\smalltext 输出四个方向的概率\\[1.5mm]
屏蔽非法动作并重新归一\\[1.5mm]
测试选择概率最大的合法动作
\end{column}
\begin{column}{.48\textwidth}
\ptext{价值分支}\\[2mm]\smalltext 估计状态价值\\[1.5mm]
与轨迹回报共同用于GAE优势估计\\[1.5mm]
剪切目标限制单次策略更新幅度
\end{column}\end{columns}\vspace{4mm}
\smalltext 训练循环：收集交互轨迹，计算优势，再更新网络。\\[2mm]
剪切系数0.2，折扣因子0.999，GAE参数0.95。\\[2mm]
当前模型的训练量和奖励设置仍需进一步检验。
\end{frame}
\note{PPO直接输出四个方向的动作概率。卷积残差网络提取棋盘特征，策略分支给出动作概率，价值分支用于估计优势。训练时先收集交互轨迹，再用GAE计算优势，并用剪切目标限制策略更新幅度，这里的剪切系数是零点二。网络在输出端屏蔽非法动作，测试时选择概率最大的合法动作。本次模型表现有限，还需要进一步检验训练量和奖励设置。}

\begin{frame}{实现与超参数}
\smalltext 继承原环境，共用配置与评测入口。\\[1mm]
\renewcommand{\arraystretch}{1.1}
\begin{tabularx}{\textwidth}{@{}l l Y@{}}\toprule
方法 & 计算预算 & 关键设置\\\midrule
Expectimax & 完整搜索与缓存 & 终局达标奖励256\\
MCTS & 每步128次模拟 & 深度12，$c=70$\\
乐观TD & 300,000局 & 初始总价值300\\
CNN-TD & 50,000局 & 学习率$3\times10^{-4}$，批量1024\\
n-tuple TD & 5,000局 & $\alpha=0.2$，探索率0.20降至0.02\\
PPO & 50,176个RL环境步 & $\gamma=0.999$，clip=0.2\\\bottomrule
\end{tabularx}\\[2mm]
搜索预算影响前瞻精度，学习率与探索率影响更新和覆盖。\\[1mm]
训练局数、RL步数和搜索次数不能直接等同，PPO另有预训练记录。
\end{frame}
\note{这些方法共用环境和评测入口。表中列出实际模型的训练量和搜索预算，不能把训练局数、环境步数和模拟次数直接等同，PPO还有教师预训练记录。超参数上，搜索预算影响前瞻精度，TD学习率和探索率影响更新与覆盖，PPO剪切系数控制更新幅度。本次没有做单因素消融，因此不能把这些值称为最优参数。}

\begin{frame}{实验结果}
\begin{columns}[T,onlytextwidth]
\begin{column}{0.60\textwidth}
\begin{tikzpicture}\begin{axis}[ybar,width=\linewidth,height=4.7cm,ymin=0,ymax=1000,
symbolic x coords={Expectimax,TD,MCTS},xtick=data,bar width=22pt,ytick={0,200,400,600,800,1000},
nodes near coords,nodes near coords style={font=\scriptsize},tick label style={font=\scriptsize},
ymajorgrids=true,grid style={gray!20},axis lines*=left,ylabel={平均终局总和}]
\addplot[fill=revisionpurple!75,draw=none] coordinates {(Expectimax,813.2666) (TD,640.743) (MCTS,282.0126)};
\end{axis}\end{tikzpicture}

\end{column}
\begin{column}{0.36\textwidth}
\smalltext\renewcommand{\arraystretch}{1.24}
\begin{tabular}{@{}lr@{}}\toprule
方法 & 达到256\\\midrule
Expectimax & 20/20\\
乐观TD & 20/20\\
CNN-TD & 16/20\\
n-tuple TD & 14/20\\
MCTS & 6/20\\
PPO & 2/20\\\bottomrule
\end{tabular}\\[3mm]
完整搜索平均总和比乐观TD高\textbf{24.3\%}。
\end{column}
\end{columns}
\smalltext 每方法20局。CNN与普通n-tuple仅有汇总值，训练与搜索预算不同。
\end{frame}
\note{结果按平均终局总和排序。完整Expectimax为761，乐观TD为612.1，两者在20局里都达到256。完整搜索比乐观TD高大约百分之二十四。CNN和普通元组表的平均总和分别为418.8和384，达标次数分别为16次和14次。当前MCTS为266，PPO为211.4。这说明完整搜索和乐观查表TD在当前设置下效果较好。但训练和搜索资源没有统一，CNN与普通元组表也只有汇总结果，所以这个榜单不能用来证明等资源下的算法优劣。}

\begin{frame}{结果解释与实验局限}
\textbf{当前设置下，完整搜索与乐观TD效果较好}\\[4mm]
\smalltext 已统一环境规则、测试种子和终止条件。\\[3mm]
训练与搜索资源尚未统一，不能据此比较等资源下的算法优劣。\\[3mm]
每方法仅20局，CNN与普通n-tuple只有汇总结果。\\[4mm]
后续固定运行时间，并用多个训练种子重复实验。
\end{frame}
\note{目前统一了环境规则、测试种子和终止条件，但计算资源尚未统一，20局也只是有限样本。下一步应固定运行时间，并用多个训练种子重复实验。}

\begingroup
\themecolor{main}
\begin{frame}{结论}
\usebeamercolor[fg]{normal text}
\large 完整Expectimax取得最高平均总和。\\[5mm]
乐观TD兼顾效果与在线计算开销。\\[8mm]
\normalsize 谢谢，欢迎提问。
\end{frame}
\endgroup
\note{本次完整Expectimax效果最好，乐观TD的在线决策开销较小。我们的汇报到这里，谢谢大家。}

\begingroup
\setbeamertemplate{footline}{\hspace{7mm}\scriptsize 补充资料\vspace{3mm}}
\begin{frame}[noframenumbering]{完整测试结果}
\smalltext\renewcommand{\arraystretch}{1.2}
\begin{tabularx}{\textwidth}{@{}Y r r r@{}}\toprule
方法 & 平均总和 & 达到256 & 平均步数\\\midrule
完整Expectimax & 761.00 & 20/20 & 344.55\\
乐观5-tuple TD & 612.10 & 20/20 & 277.60\\
CNN-TD & 418.80 & 16/20 & 188.50\\
n-tuple TD & 384.00 & 14/20 & 172.15\\
MCTS & 266.00 & 6/20 & 119.45\\
PPO & 211.40 & 2/20 & 94.55\\\bottomrule
\end{tabularx}\\[3mm]
乐观TD另一次500局测试：平均639.408，463局达标。\\[1mm]
该次测试设置不同，不与上表合并。
\end{frame}
\endgroup
\note{本页用于回答平均步数、达标次数和历史测试相关问题。500局乐观TD测试使用不同设置，平均总和639.408，463局达标，不与本页20局数据合并。}

\begingroup
\setbeamertemplate{footline}{\hspace{7mm}\scriptsize 补充资料\vspace{3mm}}
\begin{frame}[noframenumbering]{分工与参考文献}
\smalltext\renewcommand{\arraystretch}{1.15}
\begin{tabularx}{\textwidth}{@{}l Y@{}}\toprule
蔡明辰 & 环境适配、规则测试与代码整合\\
潘亮 & TD方法、模型训练与超参数整理\\
陈达 & 搜索方法、PPO与对照分析\\
孟维祎 & 评测整理、图表与报告\\\bottomrule
\end{tabularx}\\[2mm]
\fontsize{8pt}{10pt}\selectfont
[1] Guei et al. \href{https://arxiv.org/abs/2111.11090}{Optimistic Temporal Difference Learning for 2048}, 2021.\\[1mm]
[2] Szubert \& Jaskowski. N-Tuple Networks for the Game 2048, CIG 2014.\\[1mm]
[3] Schulman et al. \href{https://arxiv.org/abs/1707.06347}{Proximal Policy Optimization Algorithms}, 2017.
\end{frame}
\endgroup
\note{蔡明辰负责环境适配与代码整合，潘亮负责TD和训练，陈达负责搜索方法及PPO，孟维祎负责评测整理、图表与报告。}

\end{document}
